In \Cref{sec:game_1_problem_setting,sec:game_2_problem_setting,sec:game_3_problem_setting}, we outline the mathematical formalism for our three negotiation frameworks. In \Cref{sec:game_1_experimental_setup,sec:game_2_experimental_setup,sec:game_3_experimental_setup}, we describe the key design choices in our experimental setups, with \Cref{sec:game_1_multiagent_setup} covering our ongoing $n > 2$ extension of Game~1. Finally, \Cref{sec:cross_game_notes} collects the design notes that apply across all three games, including utility normalization, globally optimal solutions, and the solution-concept benchmarks used to measure efficiency and exploitation.

\section{Game 1: Item Allocation}

\subsection{Problem Setting}
\label{sec:game_1_problem_setting}

Agents are instructed to participate in a negotiation game.
The game consists of a set of $n$ \emph{agents} $\mathcal{N}$ presenting and voting on proposals for splitting a set of $m$ \emph{items} $\mathcal{M}$.
Each agent $i$ receives a private \emph{valuation} $\bm{v}^{i} = ( v^{i}_{1} , \dots,  v^{i}_{m} )$, where $\sum_j v^{i}_{j} = 100$ and $v^{i}_{j} \geq 0$ for all $j$, which defines how much they value each item $j$.
An \emph{allocation} $\bm{A} \in \{0,1\}^{n,m}$ is a binary matrix where $\bm{A}[i,j] = 1$ if and only if item $j$ is allocated to agent $i$.
We assume that $\sum_i \bm{A}[i,j] = 1$ for each $1 \leq j \leq m$ (no item is allocated to more than one player) and $\sum_{i,j} \bm{A}[i,j] = m$ (all items are allocated).

At each round $t$ of the negotiation game, each agent $i$ simultaneously proposes an allocation $\bm{A}^i_t$, after which all agents privately vote over all the allocations.
If an allocation $\bm{A}$ receives unanimous approval at round $t$, that allocation is selected and each agent $i$ receives \emph{utility} $u^i(\bm{A}) \coloneqq \gamma^{t-1} \bm{A}[i] \cdot \bm{v}^{i}$, where $\gamma \in (0,1]$ is a discount factor encouraging the agents to find consensus in earlier rounds and $t=1$ denotes the first round (so round~1 gives undiscounted utility).
Otherwise, the game continues to the next round.
If no allocation is selected before round $T$, the game terminates and all agents receive utility 0.

We measure the degree of competition between agents $i$ and $j$ as the \emph{cosine similarity} of their valuations, $\alpha_{ij} = \frac{{\bm{v}^i} \cdot \bm{v}^j}{\norm{\bm{v}^i} \norm{\bm{v}^j}} \in [0,1]$.
When $\alpha_{ij} = 1$, preferences are identical and every item allocated to one agent is denied to the other (pure zero-sum); when $\alpha_{ij} = 0$, valuations are orthogonal and every item can be assigned to its most-valuing agent without conflict.
For $n > 2$, we use the average pairwise similarity $\bar{\alpha} = \frac{2}{n(n-1)} \sum_{i < j} \alpha_{ij}$ as a measure of \emph{collective misalignment}.

\subsection{Experimental Setup}
\label{sec:game_1_experimental_setup}

In our current model-scale implementation, we set $n = 2$, $m = 5$, $T = 10$, and $\gamma = 0.9$.
We evaluate competition levels in $\{0.0, 0.25, 0.5, 0.75, 0.9, 0.95, 1.0\}$, number of discussion turns set at $2$, and both model orders (\texttt{weak\_first}, \texttt{strong\_first}). This yields 28 games per model pair in the base grid.
We fix GPT-5-nano as the baseline and compare against 30 adversary models.
As capability proxies, we use LMArena Elo scores \citep{chiang2024chatbotarenaopenplatform} (March 31, 2026 snapshot).

Because we instantiate the game with LLM agents, each round is structured into the following stages:

\begin{enumerate}
    \item \textbf{Initialization:} Agents receive their valuations and game instructions.
    \item \textbf{Discussion:} Agents take turns discussing their preferences, and may strategically misreport them.
    \item \textbf{Private Thinking:} Each agent consolidates takeaways in a private scratchpad.
    \item \textbf{Proposal:} Each agent submits an allocation.
    \item \textbf{Private Voting:} Agents vote yes/no on every proposal in private, so voting order cannot affect outcomes.
    \item \textbf{Evaluation:} If any proposal receives unanimous support, the negotiation terminates; ties are broken uniformly at random.
    \item \textbf{Reflection:} Agents privately reflect before the next round.
\end{enumerate}

\subsection{Multi-Agent Extension ($n > 2$)}
\label{sec:game_1_multiagent_setup}

We are currently running a multi-agent extension of Game 1 with $n \in \{3, 5\}$ agents.
This extension keeps $m = 5$, $T = 10$, $\gamma = 0.9$, and two discussion turns per round.
Competition levels are $\{0.0, 0.5, 0.9, 1.0\}$.

The extension uses paired treatment-control designs around a fixed GPT-5-nano baseline:
\begin{itemize}
    \item \textbf{Proposal 1 (invasion design):} The focal condition uses one focal model with $n-1$ GPT-5-nano agents. The matched control condition uses $n$ GPT-5-nano agents. We run 10 paired replicates for each $(\text{focal model}, n, \alpha)$ setting.
    \item \textbf{Proposal 2 (matched mixed ecology):} We sample opponent ecologies from weak, medium, and strong model pools. The focal condition replaces one baseline slot with the focal model while holding the other $n-1$ opponents fixed. The matched control condition uses GPT-5-nano in that slot with the same opponents. We run 10 matched ecologies for each $(\text{focal model}, n, \alpha)$ setting.
\end{itemize}

Each focal-control pair shares the same random seed, so causal comparisons can be made from paired deltas.
This batch is ongoing, so this subsection documents the design and not final empirical results.

\section{Game 2: Diplomatic Treaty Negotiation}

\subsection{Problem Setting}
\label{sec:game_2_problem_setting}

Agents are instructed to participate in a multi-issue diplomatic treaty negotiation game.
The game consists of a set of $n$ \emph{agents} $\mathcal{N}$ negotiating over a set of $m$ \emph{issues} $\mathcal{M}$.
An \emph{agreement} is a vector $\bm{a} = (a_1, \dots, a_m)$ where each $a_k \in [0,1]$ represents the negotiated outcome on issue $k$.
For example, if issue $k$ represents a trade tariff rate, then $a_k = 0.3$ corresponds to a 30\% tariff.

Each agent $i$ receives two private preference components: a \emph{position preference} vector $\bm{p}^{i} \in [0,1]^{m}$ specifying ideal outcomes, and an \emph{importance weight} vector $\bm{w}^{i}$ on the simplex ($w^{i}_{k} \geq 0$, $\sum_{k} w^{i}_{k} = 1$) specifying the relative salience of each issue.
Given an agreement $\bm{a}$, agent $i$'s \emph{utility} is
\begin{equation}
    u^i(\bm{a}) \coloneqq 100 \cdot \sum_{k=1}^{m} w^{i}_{k} \cdot (1 - |p^{i}_{k} - a_k|) \in [0,100].
\end{equation}
Each issue contributes utility proportional to its importance weight; the contribution is maximal at $a_k = p^{i}_{k}$ and zero when $|p^{i}_{k} - a_k| = 1$.
The factor of 100 matches the utility scale of Games 1 and 3.
Unlike the discrete allocations of Game 1, the per-issue value function $1 - |p^i_k - a_k|$ is single-peaked and continuous, so agents can make partial concessions.

The treaty game has two independent competition controls.
The first is \emph{preference correlation} $\rho \in [-1,1]$, which determines whether agents prefer similar or opposing issue values.
Position vectors are generated with a Gaussian copula~\citep{nelsen2006copulas}\footnote{A Gaussian copula is a method for generating correlated random variables with specified marginal distributions. It works by first drawing correlated normal variables, then applying the normal CDF $\Phi$ to transform each into a uniform variable, preserving the correlation structure while ensuring uniform marginals.} that preserves exact $\text{Uniform}(0,1)$ marginals:
\begin{enumerate}
    \item Sample $\bm{z} \sim \mathcal{N}(\bm{0}, \Sigma_z)$ where $(\Sigma_z)_{ij} = \rho_z$ for $i \neq j$ and $(\Sigma_z)_{ii} = 1$.
    \item Apply $p^i_k = \Phi(z^i_k)$ component-wise, where $\Phi$ is the standard normal CDF\footnote{The standard normal cumulative distribution function $\Phi(z) = \Pr(Z \leq z)$ maps any real number to $[0,1]$, converting a normal variable into a uniformly distributed one.}.
\end{enumerate}
The latent correlation $\rho_z$ is calibrated so that the Pearson correlation\footnote{The Pearson correlation coefficient measures the linear association between two random variables, ranging from $-1$ (perfect negative) to $+1$ (perfect positive).} of sampled positions equals $\rho_{\text{target}}$, using $\rho_z = 2\sin(\pi\,\rho_{\text{target}} / 6)$ \citep{kruskal1958,Embrechts_McNeil_Straumann_2002}.
For $n > 2$ agents with equicorrelation, positive semi-definiteness\footnote{A matrix is positive semi-definite (PSD) if all its eigenvalues are non-negative. Correlation matrices must be PSD; this constrains how negatively correlated multiple variables can simultaneously be.} requires $\rho_z \geq -1/(n-1)$.

The second control is \emph{interest overlap} $\theta \in [0,1]$, defined as cosine similarity between issue-importance vectors.
Weight vectors are generated with Sequential Least Squares Programming (SLSQP)~\citep{kraft1988slsqp} optimization under simplex constraints to match a target overlap.
For $n > 2$, we jointly optimize all $n$ weight vectors to match target pairwise overlap up to numerical tolerance.

Together, $\rho$ and $\theta$ separate disagreement about preferred outcomes from competition over issue salience.
Conflict is highest when preferences point in opposite directions and both agents prioritize the same issues.

At each round $t$, each agent $i$ simultaneously proposes an agreement $\bm{a}^{i}_{t}$, after which all agents privately vote over all proposals.
If a proposal $\bm{a}$ receives unanimous approval at round $t$, that agreement is selected and each agent $i$ receives utility $\gamma^{t-1} \cdot u^i(\bm{a})$, where $\gamma \in (0,1]$ is a discount factor encouraging early consensus (with $t=1$ denoting the first round).
Otherwise, the game continues to the next round.
If no agreement is reached before round $T$, the game terminates and all agents receive utility 0.

The propose-and-vote protocol places Game 2 within the multi-issue bargaining framework of \citet{rubinstein1982perfect}, extended to multiple issues in the tradition of \citet{fershtman1990importance}.

\subsection{Experimental Setup}
\label{sec:game_2_experimental_setup}

In our current conservative batch, we set $n = 2$, $m = 5$, $T = 10$, and $\gamma = 0.9$, with two discussion turns per round.
We evaluate $\rho \in \{-1.0, 0.0, 1.0\}$ and $\theta \in \{0.0, 0.5, 1.0\}$, giving 9 competition cells.
For each model pair, we run one seed per cell in both model orders, for 18 games per pair.

The play protocol follows the same seven stages as Game 1 (\Cref{sec:game_1_experimental_setup}).
The only change is at initialization: each agent receives both a position vector and an importance-weight vector, and proposes agreements in $[0,1]^m$ rather than discrete allocations.

\section{Game 3: Participatory Budgeting}

\subsection{Problem Setting}
\label{sec:game_3_problem_setting}

Agents are instructed to participate in a co-funding negotiation game.
The game consists of a set of $n$ \emph{agents} $\mathcal{N}$ negotiating over which of $m$ \emph{projects} $\mathcal{M}$ to fund by pooling their individual budgets.
Each project $j$ has a known \emph{cost} $c_j > 0$, and each agent $i$ has a private \emph{budget} $b^{i} > 0$ as well as a private \emph{valuation} $\bm{v}^{i} = (v^{i}_{1}, \dots, v^{i}_{m})$, where $\sum_j v^{i}_{j} = 100$ and $v^{i}_{j} \geq 0$ for all $j$, which defines how much agent $i$ values project $j$ being funded.
We denote the total budget $B = \sum_i b^i$ and the total project cost $C = \sum_j c_j$.

An \emph{outcome} consists of a \emph{contribution matrix} $\bm{X} \in \mathbb{R}_{\geq 0}^{n \times m}$, where $x_{ij}$ is agent $i$'s monetary contribution to project $j$.\footnote{We retain $\bm{X}$ rather than $\bm{C}$ for the contribution matrix to avoid notational collision with the total project cost $C = \sum_j c_j$.}
Each agent must respect its budget constraint: $\sum_{j} x_{ij} \leq b^i$ for all $i$.
A project $j$ is \emph{funded} if and only if it receives sufficient total contributions: $\sum_{i} x_{ij} \geq c_j$.
The \emph{funded set} is $S(\bm{X}) \coloneqq \{j \in \mathcal{M} : \sum_{i} x_{ij} \geq c_j\}$.
Contributions to unfunded projects are returned (i.e., agents pay only for projects that are successfully funded).
Given a contribution matrix $\bm{X}$, agent $i$'s \emph{utility} is
\begin{equation}
    u^i(\bm{X}) \coloneqq \sum_{j \in S(\bm{X})} v^{i}_{j} - \sum_{j \in S(\bm{X})} x_{ij}.
\end{equation}
Intuitively, agents receive the full valuation for every funded project, minus only their share of that project's cost.

Unlike Games 1 and 2, the co-funding game has threshold non-linearity: a project gives zero value until contributions meet its cost.
This places Game 3 in the class of threshold public goods games \citep{palfrey1984testing}.
The key strategic tension is coordination under free-riding pressure.

Competition is controlled by two independent parameters.
The first is \emph{preference alignment} $\alpha_{ij} = \bm{v}^{i} \cdot \bm{v}^{j} / (\norm{\bm{v}^{i}} \norm{\bm{v}^{j}}) \in [0,1]$, computed identically to Game 1 but with opposite interpretation: because the game is a public good, high $\alpha$ now signals high \emph{cooperation potential} (funded projects benefit everyone), while $\alpha = 0$ means agents value disjoint projects and co-funding provides no advantage.
Valuation vectors with target alignment are generated by SLSQP minimization of $(\cos(\bm{v}^1, \bm{v}^2) - \alpha_{\text{target}})^2$ subject to $v^i_j \geq 0$ and $\sum_j v^i_j = 100$, the same routine used for Game 2's interest overlap; for $n > 2$ we report $\bar{\alpha} = \frac{2}{n(n-1)} \sum_{i < j} \alpha_{ij}$.

The second parameter is \emph{budget scarcity} $\sigma \in (0,1]$, defined as the ratio of total budget to total project cost:
\begin{equation}
    \sigma = \frac{\sum_{i} b^i}{\sum_{j} c_j} = \frac{B}{C}.
\end{equation}
When $\sigma = 1$, the agents can collectively afford all projects and competition is minimal.
As $\sigma \to 0$, agents can collectively afford almost nothing and every dollar is contested.
When per-agent budgets are equal ($b^i = B/n$) and $b^i < \min_j c_j$, no single agent can fund even the cheapest project alone, forcing agents into coalition dynamics for \emph{every} project.

The two parameters are orthogonal: $(\alpha, \sigma) = (1, 0.2)$ is a high-alignment but scarce-budget regime, while $(0, 1)$ is a low-alignment but resource-rich one.

Game 3 uses the same propose-and-vote loop as Games 1 and 2, with one adaptation for the public-good setting: each agent submits a contribution vector over projects, the environment aggregates these into a single joint proposal, computes the funded set under the threshold constraint, and then all agents privately vote accept/reject on that joint proposal.
Unanimous acceptance in round $t$ locks in the proposal with utility discounted by $\gamma^{t-1}$; otherwise play advances to round $t+1$, and if no agreement is reached by round $T$ all agents receive utility 0.

\subsection{Experimental Setup}
\label{sec:game_3_experimental_setup}

In our current conservative batch, we set $n = 2$, $m = 5$, $T = 10$, and $\gamma = 0.9$.
We evaluate $\alpha \in \{0.0, 0.5, 1.0\}$ and $\sigma \in \{0.2, 0.6, 1.0\}$, yielding 9 competition cells.
Project costs are drawn uniformly from $[10, 30]$.
Given sampled costs, budgets are set symmetrically across agents so that total budget is approximately $\sigma C$ after integer rounding in the environment.
For each model pair, we run one seed per cell in both model orders, for 18 games per pair.
Discussion turns are fixed at 2 with own-transparency feedback (aggregate project totals are visible; per-agent contributions are private).

Each round proceeds through:

\begin{enumerate}
    \item \textbf{Initialization:} Each agent receives its valuation vector, its budget, and the costs of all projects.
    \item \textbf{Communication:} Agents take turns sending natural-language messages about priorities and possible coalitions.
    \item \textbf{Private Thinking:} Each agent updates strategy in a private scratchpad.
    \item \textbf{Contribution Proposal:} Each agent $i$ submits $\bm{x}^i_t = (x^i_{t,1}, \ldots, x^i_{t,m})$ with $\sum_j x^i_{t,j} \leq b^i$.
    \item \textbf{Aggregation:} The environment aggregates submitted vectors into one joint proposal and computes the implied funded set.
    \item \textbf{Private Voting:} Agents vote accept or reject on the joint proposal.
    \item \textbf{Evaluation and Reflection:} If acceptance is unanimous, the game ends with that proposal. Otherwise, the next round begins.
\end{enumerate}

If no unanimous acceptance occurs by round $T$, the game ends in failure with utility 0 for all agents.

\subsection{Behavioral Metrics and Scalable Evaluation}
\label{sec:game_3_behavioral_metrics}

The Game 3 protocol generates two parallel data streams per game: (1) the sequence of contribution vectors $\{\bm{x}^i_t\}_{i,t}$, which are structured numerical data, and (2) the natural-language transcripts from communication phases.
We exploit this structure to define behavioral metrics that are computable at scale across thousands of games.

\paragraph{Efficiency and individual metrics.}
We measure \emph{utilitarian efficiency} $\eta = \\ SW_{\text{actual}} / SW_{\text{optimal}}$, where $SW_{\text{optimal}}$ is the social welfare of the welfare-maximizing funded set (a knapsack problem).
At the individual level, we compute each agent's utility $u^i$ and a \emph{free-rider index} $F_{ij} = (v^i_j / \sum_k v^k_j) / (x_{ij} / \sum_k x_{kj})$ for each funded project $j$ to which agent $i$ contributed ($x_{ij} > 0$), where $F_{ij} > 1$ indicates that agent $i$ extracts disproportionate value relative to their cost share.
Agent-project pairs where $x_{ij} = 0$ are classified as \emph{pure free-riding} and tracked separately.
We also define an \emph{exploitation index} $E_i = (u^i_{\text{actual}} - u^i_{\text{Lindahl}}) / |u^i_{\text{Lindahl}}|$ analogous to Game 2's exploitation index relative to the NBS: when $E_i > 0$, agent $i$ extracted more utility than they would receive under the proportionally fair Lindahl cost-sharing (defined in \Cref{sec:solution_concepts}), enabling analysis of whether the Elo gap predicts asymmetric exploitation.
We also report the \emph{underfunding rate}: the fraction of \emph{budget-feasible} positive-surplus projects that go unfunded.
Specifically, a project $j$ contributes to the underfunding rate only if $\sum_i v^i_j > c_j$ (positive surplus) \emph{and} there exists a feasible reallocation of the agents' budgets that could fund $j$ alongside the actually funded projects.
This conditions on coordination being \emph{in principle possible}, distinguishing genuine coordination failure (agents could have funded it but didn't) from budget impossibility (there simply weren't enough resources).

\paragraph{Strategic behavior metrics.}
Four additional metrics link communication content to numerical behavior:
\begin{itemize}
    \item \textbf{Promise-keeping rate}: the fraction of extracted contribution commitments (e.g., ``I will contribute 15 to Project 2'') that match the agent's subsequent pledge within a tolerance, measuring consistency versus deception.
    \item \textbf{Adaptation rate}: $\text{Adapt}_i = \frac{1}{(T-1)\,b^i} \sum_{t=2}^{T} \|\bm{x}^i_t - \bm{x}^i_{t-1}\|_1$, computable purely from contribution vectors.
    \item \textbf{Persuasion effectiveness}: a Granger-style measure of whether agent $i$ advocating for project $j$ at round $t$ predicts an increase in other agents' contributions to $j$ at round $t+1$, with positive and negative advocacy treated as separate signals.
    \item \textbf{Coalition formation}: detected as 2+ agents repeatedly co-funding the same project(s) across $\lceil T/2 \rceil$ or more consecutive rounds.
\end{itemize}

\paragraph{Scalable extraction pipeline.}
To compute behavioral metrics from thousands of game transcripts across the active model roster, we use a three-stage pipeline.
First, a lightweight model (Haiku-class) processes each round's communication phase and produces a per-round structured extraction: commitments (agent, project, amount), advocacy events (agent, project, sentiment), and threats.
The sentiment field in advocacy events captures both the direction (for or against) and intensity of each mention, addressing the distinction between positive and negative advocacy.
Second, all metrics are computed from the structured data (contribution vectors combined with extracted events); promise-keeping, adaptation, coalition detection, and persuasion are all reducible to numerical comparisons once communication events are extracted.
Third, a random sample of transcripts is reviewed by an LLM judge operating under a detailed rubric to validate the extraction pipeline's accuracy.
This approach scales linearly in cost with the number of games while remaining auditable through spot-checks on the extraction stage.

\section{Cross-Game Design Notes}
\label{sec:cross_game_notes}

\subsection{Utility Normalization Across Games}
\label{sec:normalization}

All three games are designed so that raw utilities share a $[0, 100]$ scale, but the \emph{achievable} maximum differs across settings and must be handled carefully.

In Games 1 and 2, the simplex constraint on valuations ($\sum_j v^i_j = 100$) and on importance weights ($\sum_k w^i_k = 1$), combined with a per-issue value function in $[0,1]$ for Game 2, guarantees $u^i \in [0, 100]$: an agent achieves 100 by receiving every item (Game 1) or by the agreement matching their ideal position on every issue (Game 2), and 0 at timeout.
Raw utilities are therefore directly comparable across $\alpha$, $\rho$, and $\theta$ levels.
Position vectors $\bm{p}^i$ in Game 2 are intentionally not on a simplex: positions are independent ideals per issue, so a sum constraint would be semantically inappropriate.

In Game 3, the simplex constraint still bounds $u^i \leq 100$ in principle (pure free-riding on all funded projects), but the \emph{achievable} maximum shrinks with budget scarcity $\sigma$: at low $\sigma$ even perfect cooperation cannot fund every project.
Cross-$\sigma$ comparisons therefore use the efficiency metric $\eta = SW_{\text{actual}} / SW_{\text{optimal}}$, which normalizes by the knapsack optimum for the given $\sigma$.
Raw utilities remain comparable across $\alpha$ at fixed $\sigma$.

\subsection{Computing the Globally Optimal Solution}
\label{sec:optimal_computation}

For all three games, the socially optimal outcome---and thus the denominator of $\eta = SW_{\text{actual}} / SW_{\text{optimal}}$---can be computed exactly at our experimental scale ($n = 2$, $m = 5$).

\paragraph{Game 1.}
The utilitarian optimum assigns each item $j$ to the agent $i^* = \arg\max_i v^i_j$ who values it most.
Since items are rivalrous but independent (allocating item $j$ to agent $i$ does not affect the value of allocating item $j'$ to any agent), greedy assignment is globally optimal.
This runs in $O(nm)$, which is 10 comparisons for $n=2$, $m=5$.

\paragraph{Game 2.}
The social welfare function decomposes per issue: $SW(\bm{a}) = \sum_k \sum_i w^i_k (1 - |p^i_k - a_k|)$.
Since each issue $k$ contributes independently, the optimum sets each $a_k$ to the \emph{weighted median} of agent positions $\{p^i_k\}$ with weights $\{w^i_k\}$, which minimizes the weighted sum of absolute deviations.
For $n = 2$, the optimal $a_k$ is simply the ideal position of whichever agent has higher weight on issue $k$.
This runs in $O(mn \log n)$.

\paragraph{Game 3.}
The socially optimal funded set $S^*$ maximizes $\sum_{j \in S} (\sum_i v^i_j - c_j)$ subject to $\sum_{j \in S} c_j \leq B$.
This is a 0-1 knapsack problem with ``profit'' $\pi_j = \sum_i v^i_j - c_j$ and ``weight'' $c_j$ for each project $j$.
While knapsack is NP-hard in general, for $m = 5$ projects we enumerate all $2^5 = 32$ subsets, check budget feasibility, and select the maximum-surplus set.
For larger $m$, the standard pseudo-polynomial dynamic programming algorithm runs in $O(m \cdot B)$.
Note that the knapsack gives the optimal \emph{funded set} but not the optimal cost-sharing: social welfare is independent of how costs are split (individual contributions cancel in the sum), so $SW^*$ depends only on $S^*$.

\subsection{Solution Concept Benchmarks}
\label{sec:solution_concepts}

The social optimum identifies the best \emph{total} outcome but does not specify a fair split, and the Nash equilibria of the extensive-form game are massively multiple (including degenerate equilibria where every agent rejects every proposal).
We therefore use cooperative-game-theoretic benchmarks, which give unique and axiomatically justified predictions, and report each game's outcome as a distance from its benchmark.

\paragraph{Games 1 and 2: Nash Bargaining Solution (NBS).}
The NBS maximizes the product of agents' gains over the disagreement point (utility 0 from timeout):
\begin{equation}
    \bm{A}^{\text{NBS}} = \arg\max_{\bm{A}} \prod_{i=1}^{n} u^i(\bm{A})
\end{equation}
The NBS is unique, Pareto efficient, and captures the notion of a fair outcome under equal bargaining power.
It balances efficiency (high total utility) with equity (neither agent receives too little).

For Game 1, we enumerate all $2^m = 32$ allocations and select the one maximizing $u^1(\bm{A}) \cdot u^2(\bm{A})$.
For Game 2, the Nash product over $\bm{a} \in [0,1]^m$ is smooth, so we solve it with L-BFGS-B (Limited-Memory Bounded Broyden-Fletcher-Goldfarb-Shanno optimization).

We measure \emph{distance from NBS}:
\begin{equation}
    D_{\text{NBS}} = \left\|\bm{u}_{\text{actual}} - \bm{u}_{\text{NBS}}\right\|_2
\end{equation}
and the \emph{exploitation index} for each agent $i$:
\begin{equation}
    E_i = \frac{u^i_{\text{actual}} - u^i_{\text{NBS}}}{u^i_{\text{NBS}}}
\end{equation}
When $E_i > 0$, agent $i$ extracted more than their fair share; when $E_i < 0$, they were exploited.
This enables analysis of whether stronger models systematically exploit weaker negotiation partners: if the Elo gap between models predicts asymmetric $E_i$ values, that is evidence of strategic exploitation scaling with capability.

\paragraph{Game 3: Lindahl equilibrium.}
In the co-funding game, the NBS is less natural because agents submit \emph{individual} contribution vectors (no joint proposal to optimize over).
Instead, the \emph{fairness} benchmark is the \emph{Lindahl equilibrium}, in which each agent pays for each funded project in proportion to their valuation:
\begin{equation}
    x^{\text{Lindahl}}_{ij} = c_j \cdot \frac{v^i_j}{\sum_k v^k_j} \quad \text{for each funded project } j \in S^*
\end{equation}
Note that the Lindahl equilibrium addresses the fairness question (``who pays how much?'') rather than the efficiency question (``which projects are funded?''), since efficiency depends only on the funded set $S^*$ and is measured separately by $\eta$.

The \emph{free-rider index} $F_{ij}$ is defined as the ratio of agent $i$'s value share to their cost share for funded project $j$:
\begin{equation}
    F_{ij} = \frac{v^i_j / \sum_k v^k_j}{x_{ij} / \sum_k x_{kj}} \quad \text{for } j \in S, \; x_{ij} > 0.
\end{equation}
Intuitively, if agent $i$ receives 40\% of the total value from project $j$ but pays only 20\% of the cost, then $F_{ij} = 2$: they are extracting twice their ``fair share.''
At the Lindahl equilibrium, $F_{ij} = 1$ for every agent on every project, so nobody overpays or underpays relative to their benefit.

We measure \emph{distance from Lindahl}:
\begin{equation}
    D_{\text{Lindahl}} = \left\|\bm{X}_{\text{actual}} - \bm{X}_{\text{Lindahl}}\right\|_F
\end{equation}
where $\|\cdot\|_F$ is the Frobenius norm (root sum of squared entry-wise differences).
Note that $D_{\text{NBS}}$ and $D_{\text{Lindahl}}$ are not directly comparable across games because they measure distances in different spaces (utility space for $D_{\text{NBS}}$ vs.\ contribution-matrix space for $D_{\text{Lindahl}}$).
However, the exploitation indices $E_i$ are comparable: in both cases, $E_i > 0$ means agent $i$ got more than the fair benchmark, and $E_i < 0$ means they got less.

% \paragraph{Game 3: Core of the cooperative game.}
% The \emph{core} is the set of cost-sharing arrangements where no coalition of agents can do better by breaking away and funding projects independently.
% Formally, a cost-sharing $\bm{X}$ for funded set $S$ is in the core if and only if for every coalition $T \subseteq \mathcal{N}$:
% \begin{equation}
%     \sum_{i \in T} u^i(\bm{X}, S) \geq V(T)
% \end{equation}
% where $V(T)$ is the value of coalition $T$, meaning the maximum total utility agents in $T$ could achieve using only their own budgets:
% \begin{equation}
%     V(T) = \max_{S' \subseteq \mathcal{M}} \sum_{j \in S'} \left(\sum_{i \in T} v^i_j - c_j\right) \quad \text{s.t.} \quad \sum_{j \in S'} c_j \leq \sum_{i \in T} b^i
% \end{equation}
% Computing $V(T)$ is a knapsack problem; for $m=5$, each coalition's value is found by enumerating 32 subsets.
% With $n=2$ agents, there are $2^n - 1 = 3$ non-empty coalitions to check.
% With $n=5$, there are 31 coalitions, each requiring 32 subset evaluations (under 1000 checks total).

% The core is a polytope in the space of contribution matrices.
% We report whether the actual cost-sharing lies inside the core and, if not, its minimum distance to the core boundary (computable via linear programming).

% \paragraph{Efficiency-fairness decomposition.}
% For all three games, deviation from the optimum decomposes into two orthogonal components: an \emph{efficiency loss} $SW^* - SW_{\text{actual}}$ (surplus destroyed by wrong allocations, unfunded projects, or timeout) and a \emph{fairness deviation} measured by distance from NBS (Games 1, 2) or Lindahl (Game 3) among outcomes with the same $SW_{\text{actual}}$.
% This separates ``they left money on the table'' from ``they captured the right total but split it unevenly,'' sharpening the analysis of how capability affects negotiation outcomes.
% For Game 3, the core analysis adds a stronger stability check: whether a coalition could have strictly improved by breaking away from the grand coalition.
